% =======================================================================
% ΚΡΙΤΗΡΙΟ ΑΞΙΟΛΟΓΗΣΗΣ — ΔΟΚΙΜΑΣΤΙΚΟ ΠΕΡΙΕΧΟΜΕΝΟ
% =======================================================================
% Τα συγκεκριμένα διαγωνίσματα αφορούν την Ενότητα 1 (Κινηματική).
% Τοποθετούνται πριν από το επόμενο κεφάλαιο στο book-content.tex,
% ώστε να εμφανίζονται ΕΝΤΟΣ της Ενότητας 1 στα περιεχόμενα.
% Ολες οι μορφοποιήσεις ορίζονται κεντρικά.
% Τα παρακάτω ΔΕΝ αποτελούν ολοκληρωμένο επιστημονικά διαγώνισμα προς χρήση.
% Είναι σχεδιαστικές δοκιμές των περιβαλλόντων και των σημείων.
\begin{Assessment}[student=true,date=true,duration={90 λεπτά}]{Ευθύγραμμη κίνηση}
\AssessmentTopic{Α}
\begin{AssessmentQuestion}[5]\label{asq:KRI1-A1}
Η ταχύτητα ενός σώματος παραμένει σταθερή κατά μέτρο και κατεύθυνση.
Η επιτάχυνσή του είναι:
\begin{enumerate}[label=\BookGreek*. ,leftmargin=7mm,itemsep=1pt]
\item θετική. \item αρνητική. \item μηδενική. \item μεταβαλλόμενη.
\end{enumerate}
\end{AssessmentQuestion}
\begin{AssessmentQuestion}[5]
Να χαρακτηρίσετε την πρόταση ως σωστή ή λανθασμένη:
«Το διάστημα είναι πάντοτε ίσο με το μέτρο της μετατόπισης».
\TrueFalseChoices
\end{AssessmentQuestion}
\begin{AssessmentQuestion}[5]
Το μέτρο της μέσης διανυσματικής ταχύτητας ενός σώματος
υπολογίζεται από τη σχέση $|\bar{\upsilon}|=\frac{|\Delta x|}{\Delta t}$.
Να αναφέρετε τις μονάδες SI των μεγεθών $\Delta x$ και $\Delta t$.
\end{AssessmentQuestion}
\begin{AssessmentQuestion}[5]
Ένα κινητό κινείται από τη θέση $x_1=-2\,\mathrm{m}$ στη θέση
$x_2=6\,\mathrm{m}$. Ποια είναι η αλγεβρική μετατόπισή του;
\end{AssessmentQuestion}
\begin{AssessmentQuestion}[5]
Να ορίσετε τη μέση αριθμητική ταχύτητα σε μια κίνηση.
\end{AssessmentQuestion}

\AssessmentTopic{Β}
\begin{AssessmentQuestion}[12]
Δύο κινητά Α και Β κινούνται στον ίδιο άξονα.
Τα διαγράμματα θέσης--χρόνου φαίνονται παρακάτω.
\begin{center}
\begin{tikzpicture}[scale=.67]
 \draw[book/axis] (0,0)--(6.4,0) node[right]{$t\ (\mathrm{s})$};
 \draw[book/axis] (0,-.4)--(0,3.9) node[above]{$x\ (\mathrm{m})$};
 \draw[book/velocity] (0,.35)--(5.4,3.55) node[right]{A};
 \draw[book/accel] (0,3.1)--(5.4,1.05) node[right]{B};
\end{tikzpicture}
\end{center}
Ποιο κινητό έχει θετική και ποιο αρνητική αλγεβρική ταχύτητα;
Να αιτιολογήσετε την επιλογή σας.
\end{AssessmentQuestion}
\begin{AssessmentQuestion}[13]
Δύο σώματα διανύουν σε ίσους χρόνους διαστήματα $12\,\mathrm{m}$ και
$18\,\mathrm{m}$ αντίστοιχα. Ποιο έχει μεγαλύτερη μέση αριθμητική ταχύτητα;
\end{AssessmentQuestion}

\AssessmentTopic{Γ}
\begin{AssessmentQuestion}[10]\label{asq:KRI1-G1}
Ένα κινητό ξεκινά από τη θέση $x_0=2\,\mathrm{m}$ και κινείται με
σταθερή ταχύτητα $\upsilon=3\,\mathrm{m/s}$.
Να γράψετε την εξίσωση της θέσης του.
\end{AssessmentQuestion}
\begin{AssessmentQuestion}[15]
Για το προηγούμενο κινητό, να υπολογίσετε τη θέση του τη χρονική στιγμή
$t=4\,\mathrm{s}$ και τη μετατόπισή του στο ίδιο χρονικό διάστημα.
\end{AssessmentQuestion}

\AssessmentTopic{Δ}
\begin{AssessmentQuestion}[9]
Ένα σώμα αρχίζει να κινείται από την ηρεμία με σταθερή επιτάχυνση
$\alpha=2\,\mathrm{m/s^2}$. Να υπολογίσετε την ταχύτητά του τη στιγμή
$t=5\,\mathrm{s}$.
\end{AssessmentQuestion}
\begin{AssessmentQuestion}[8]
Να υπολογίσετε τη μετατόπισή του στα πρώτα $5\,\mathrm{s}$.
\end{AssessmentQuestion}
\begin{AssessmentQuestion}[8]
Να σχεδιάσετε ποιοτικά το διάγραμμα ταχύτητας--χρόνου για
$0\leq t\leq5\,\mathrm{s}$.
\end{AssessmentQuestion}
% Προαιρετικός έλεγχος στο αρχείο LOG: 25+25+25+25=100.
\AssessmentPointsCheck[100]
\end{Assessment}

% Δεύτερο, σύντομο κριτήριο: ΟΥΤΕ στοιχεία μαθητή ΟΥΤΕ ημερομηνία
% ΟΥΤΕ διάρκεια. Κάποια ερωτήματα δεν έχουν μονάδες.
\begin{Assessment}{Σύντομη επανάληψη: μετατόπιση}
\AssessmentTopic{Α}
\begin{AssessmentQuestion}
Ποια είναι η διαφορά ανάμεσα στη μετατόπιση και στο διάστημα;
\end{AssessmentQuestion}
\begin{AssessmentQuestion}[5]
Δίνεται $x_A=4\,\mathrm{m}$ και $x_B=-1\,\mathrm{m}$.
Να βρείτε τη μετατόπιση από το Α στο Β.
\end{AssessmentQuestion}
\AssessmentTopic{Β}
\begin{AssessmentQuestion}
Να εξηγήσετε αν είναι δυνατόν η μετατόπιση ενός κινητού να είναι μηδέν,
ενώ το διάστημα που έχει διανύσει είναι μεγαλύτερο από μηδέν.
\end{AssessmentQuestion}
\end{Assessment}
